\section{Tugas dan Ruang Lingkup Analisis Semantik}
\label{sec:ch14-ruang}

\subsection{Di Luar Sintaksis}

Analisis sintaksis memastikan program memiliki bentuk yang benar.
Analisis semantik memastikan program memiliki makna yang masuk akal.
Kesalahan seperti variabel yang tidak dideklarasikan, tipe yang tidak
cocok, atau pemanggilan fungsi dengan jumlah argumen yang salah tidak
dapat dideteksi oleh tata bahasa bebas konteks.

\begin{table}[H]
\centering
\caption{Tugas utama analisis semantik}
\label{tab:ch14-tugas}
\begin{tabularx}{\textwidth}{@{}l X@{}}
\toprule
\textbf{Tugas} & \textbf{Pertanyaan yang dijawab} \\
\midrule
Resolusi nama & Apakah identifier sudah dideklarasikan? \\
Pemeriksaan tipe & Apakah operand memiliki tipe yang kompatibel? \\
Pemeriksaan fungsi & Apakah jumlah dan tipe argumen sesuai? \\
Kontrol alur & Apakah \code{break} berada di dalam loop? \\
Inisialisasi & Apakah variabel dipakai sebelum diberi nilai? \\
Keterjangkauan & Apakah semua jalur fungsi mengembalikan nilai? \\
\bottomrule
\end{tabularx}
\end{table}

\subsection{Informasi yang Dibutuhkan}

Analisis semantik membutuhkan:
\begin{itemize}[leftmargin=*]
  \item AST atau pohon urai yang telah dibangun;
  \item tabel simbol dengan informasi scope dan tipe;
  \item aturan tipe untuk setiap operator dan konstruksi;
  \item mekanisme pelaporan kesalahan.
\end{itemize}

Tanpa tabel simbol, analisis hanya dapat memeriksa struktur lokal.
Tanpa aturan tipe, analisis tidak dapat menentukan kompatibilitas
operand.

\subsection{Urutan Analisis}

Urutan yang umum:
\begin{enumerate}[label=\arabic*.,leftmargin=*]
  \item bangun tabel simbol dan resolusi nama;
  \item tentukan tipe setiap ekspresi;
  \item periksa kecocokan tipe pada assignment dan pemanggilan;
  \item periksa aturan kontrol alur;
  \item laporkan kesalahan dan lanjutkan.
\end{enumerate}

Analisis dapat dilakukan dalam satu lintasan AST atau beberapa
lintasan. Beberapa lintasan memudahkan pemisahan tugas, tetapi
memerlukan AST yang stabil.

\begin{infobox}
Analisis semantik tidak harus menghasilkan kode. Tujuannya adalah
memastikan program valid dan mengumpulkan informasi yang dibutuhkan
fase berikutnya, seperti tipe dan alamat simbol.
\end{infobox}
